\documentclass[thesis/Elad_Kliger_Thesis.tex]{subfiles}

\makeatletter
\@ifundefined{preamble@loaded}{
    \def\preamble@loaded{}

\makeatletter

\usepackage[utf8]{inputenc}

\usepackage{adjustbox}
\usepackage{algpseudocode}
\usepackage{amsfonts}
\usepackage{amsmath}
\usepackage{amssymb}
\usepackage{amsthm}
\usepackage{array}
% \usepackage[noaffiliation]{authblk}
\usepackage[english]{babel}
\usepackage{bm}
\usepackage{color}
\usepackage{csquotes}
\usepackage{dsfont}
\usepackage{etoolbox}
\usepackage{faktor}
\usepackage[T1]{fontenc}
\usepackage{graphicx}
\usepackage[acronym,toc]{glossaries}
% Keep the paper on REVTeX's native page grid and float dimensions.
\ifdefined\paper\else
    \usepackage{geometry}
\fi
\usepackage{mathtools}
\usepackage{microtype} % improves spacing overall
\usepackage{physics}
\usepackage{soul}
\usepackage{tikz}
\usepackage{tikz-cd}

\@ifclassloaded{beamer}{  % Context dependent inclusions
    \usepackage{algorithm}
    % Computer Modern Sans names these available shapes differently.
    \DeclareFontShape{T1}{cmss}{b}{n}{<->ssub*cmss/bx/n}{}
    \DeclareFontShape{T1}{cmss}{m}{sc}{<->ssub*cmr/m/sc}{}
}{
    \ifdefined\thesis
        \geometry{margin=1in}
    \fi
}

\usetikzlibrary{arrows}
\tikzcdset{arrow style=tikz,
diagrams={>={Stealth[round,length=4pt,width=6pt,inset=3pt]}}}

\setlength{\marginparwidth}{2cm}
\usepackage{todonotes}

\usepackage{xcolor}
\usepackage{xfrac}

\usepackage{subfiles} % Best loaded last

\makeatother

% =============================================================================
% COMMON DEFINITIONS AND SHORTCUTS
% =============================================================================
% This file contains LaTeX function definitions and shortcuts to be included
% in other tex files. Include with: % =============================================================================
% COMMON DEFINITIONS AND SHORTCUTS
% =============================================================================
% This file contains LaTeX function definitions and shortcuts to be included
% in other tex files. Include with: \input{common/headers/definitions.tex}


% =============================================================================
% APS recommendattions
% =============================================================================
% \mathfrak{#1} % Fraktur
% \mathbb{#1}   % Blackboard

% single line eq:
% (\begin{widetext})      % protect from two column mode in release mode
% (\begin{subequations})  % for multiple equations with a shared number
% \begin{eqnarray(*)}
% ...
% (\nonumber\\)
% \label{eq:one}.
% \end{eqnarray(*)}
% (\end{subequations})

% unnumbered equations
% \[
% ...
% \]

% figures should be put into the text as floats.
% surround with turnpage environment for landscape figure

% \begin{figure}      % figure* for full page width
% \label{}
% \includegraphics[width=\textwidth, trim=0 10bp 0 0, clip]{}
% \caption{}
% \end{figure}

% tables should appear as floats within the text

% \begin{table}%[H] add [H] placement to break table across pages
% \caption{\label{}}
% \begin{ruledtabular}
% \begin{tabular}{}
% Lines of table here ending with \\
% \end{tabular}
% \end{ruledtabular}
% \end{table}


% =============================================================================
% NOTATIONS
% =============================================================================

%% Data sets
\newcommand{\eventsymbol}{{x}}
\newcommand{\anyevent}{\vec{\eventsymbol}}
\newcommand{\event}{{\vec{\eventsymbol_i}}}
\newcommand{\dataset}[1]{{\mathbf{#1}}}         % A dataset
\newcommand{\true}[1]{{\dataset{#1}}_{\text{true}}}  % A true dataset
\newcommand{\crdataset}[1]{{\dataset{#1}^{CR}}}
\newcommand{\srdataset}[1]{{\dataset{#1}^{SR}}}

%% Detecting events
\newcommand{\detector}{\delta}            % Detector function

%% Hypothesis testing
\newcommand{\hypothesis}[1]{{\mathcal{H}_{#1}}}
\newcommand{\nullhypothesis}{{\hypothesis{0}}}
%%% Parameterizations and pdfs
\newcommand{\pdf}{{\lambda}}
\newcommand{\param}{{\mu}}
\newcommand{\nparam}{{\theta}}
\newcommand{\parametrized}[1]{{#1}_{\nparam}^{\param}}
\newcommand{\parampdf}{{\parametrized{\pdf}}}
\newcommand{\paramf}{{f^{\param}}}

%% Hypotheses
%%% Distributions
\newcommand{\dist}[1]{{\pdf_{#1}(\anyevent)}}   % Any distribution of x
\newcommand{\srdist}[1]{{\pdf_{#1}^{SR}(\anyevent)}}   % Any distribution of x
\newcommand{\crdist}[1]{{\pdf_{#1}^{CR}(\anyevent)}}   % Any distribution of x
\newcommand{\distevent}[1]{\pdf_{#1}(\event)}
\newcommand{\srdistevent}[1]{\pdf_{#1}^{SR}(\event)}
\newcommand{\crdistevent}[1]{\pdf_{#1}^{CR}(\event)}
\newcommand{\refdist}[1][\anyevent]{n_{\mathcal{R}}({#1})}  % Reference distribution
\newcommand{\crrefdist}[1][\anyevent]{n_{\mathcal{R}}^{CR}({#1})}
\newcommand{\srrefdist}[1][\anyevent]{n_{\mathcal{R}}^{SR}({#1})}

%% Likelihood
\newcommand{\likelihood}{\mathcal{L}}           % Likelihood function
\newcommand{\llikelihood}{\log\mathcal{L}}      % Log-likelihood function
\newcommand{\plratio}{\Lambda(\param)}
\newcommand{\invllratio}{{q_{0}}}

% =============================================================================
% Theorems
% =============================================================================

\newtheorem{theorem}{Theorem}
\newtheorem{definition}[theorem]{Definition}
\newtheorem{conjecture}[theorem]{Conjecture}
\newtheorem{lemma}[theorem]{Lemma}
\newtheorem{notation}[theorem]{Notation}
\newtheorem{observation}[theorem]{Observation}
\newtheorem{proposition}[theorem]{Proposition}
\newtheorem{property}[theorem]{Property}
\newtheorem{example}[theorem]{Example}
\newtheorem{remark}[theorem]{Remark}
\newtheorem{corollary}[theorem]{Corollary}

\newtheorem*{theorem*}{Theorem}

% =============================================================================
% Display Macros
% =============================================================================

% Size relative to the enclosing column, minipage, or slide column.
\newcommand{\addsizedgraph}[2]{%
        \centering
        \includegraphics[
            width=#1\linewidth,
            keepaspectratio,
            trim=0 24bp 0 0,
            clip
        ]{#2}
}
\newcommand{\addgraph}[1]{\addsizedgraph{1}{#1}}

\newcommand{\addsizedpic}[2]{%
        \centering
        \includegraphics[
            width=#1\linewidth,
            keepaspectratio
        ]{#2}
}
\newcommand{\addpic}[1]{\addsizedpic{1}{#1}}

% Visible draft replacement for figures whose assets are not available yet.
% Unlike \includegraphics{}, this keeps incomplete manuscripts compilable.
\newcommand{\placeholdergraphic}[1][Figure pending]{%
    \centering
    \fbox{\parbox[c][0.18\textheight][c]{0.85\linewidth}{%
        \centering\textit{#1}%
    }}%
}

% Macro to insert a full-page PDF without altering document orientation
\newcommand{\includefullpage}[2][1]{%
  \clearpage
  \thispagestyle{empty}% Suppress header/footer for the inserted page
  \noindent
  \begin{picture}(0,0)
    \put(-72,-750){% Adjust offset if needed for your margins
      \includegraphics[page=#1,width=\paperwidth,height=\paperheight,keepaspectratio=false]{#2}%
    }
  \end{picture}%
  \clearpage
}

% =============================================================================
% DRAFT MARKERS (TODO/FIXME/TBD)
% =============================================================================
% Usage:
%   \TODO{finish this paragraph}
%   \FIXME{replace placeholder with measured value}
%   \TBD{insert citation}
%
% Release behavior:
%   - Default is draft mode (markers visible in PDF).
%   - Set \draftfalse (after loading definitions) to make unresolved markers
%     fail the build.
\makeatletter
\@ifundefined{ifdraft}{\newif\ifdraft\drafttrue}{}
\makeatother

\ifdraft
  \newcommand{\draftmarker}[2]{%
    \textcolor{#1}{\textbf{[#2]}}%
  }
  \newcommand{\TODO}[1]{\draftmarker{red}{TODO: #1}}
  \newcommand{\FIXME}[1]{\draftmarker{orange}{FIXME: #1}}
  \newcommand{\TBD}[1]{\draftmarker{blue}{TBD: #1}}
\else
  \newcommand{\TODO}[1]{\errmessage{Unresolved TODO: #1}}
  \newcommand{\FIXME}[1]{\errmessage{Unresolved FIXME: #1}}
  \newcommand{\TBD}[1]{\errmessage{Unresolved TBD: #1}}
\fi

% =============================================================================
% END OF DEFINITIONS
% =============================================================================



% =============================================================================
% APS recommendattions
% =============================================================================
% \mathfrak{#1} % Fraktur
% \mathbb{#1}   % Blackboard

% single line eq:
% (\begin{widetext})      % protect from two column mode in release mode
% (\begin{subequations})  % for multiple equations with a shared number
% \begin{eqnarray(*)}
% ...
% (\nonumber\\)
% \label{eq:one}.
% \end{eqnarray(*)}
% (\end{subequations})

% unnumbered equations
% \[
% ...
% \]

% figures should be put into the text as floats.
% surround with turnpage environment for landscape figure

% \begin{figure}      % figure* for full page width
% \label{}
% \includegraphics[width=\textwidth, trim=0 10bp 0 0, clip]{}
% \caption{}
% \end{figure}

% tables should appear as floats within the text

% \begin{table}%[H] add [H] placement to break table across pages
% \caption{\label{}}
% \begin{ruledtabular}
% \begin{tabular}{}
% Lines of table here ending with \\
% \end{tabular}
% \end{ruledtabular}
% \end{table}


% =============================================================================
% NOTATIONS
% =============================================================================

%% Data sets
\newcommand{\eventsymbol}{{x}}
\newcommand{\anyevent}{\vec{\eventsymbol}}
\newcommand{\event}{{\vec{\eventsymbol_i}}}
\newcommand{\dataset}[1]{{\mathbf{#1}}}         % A dataset
\newcommand{\true}[1]{{\dataset{#1}}_{\text{true}}}  % A true dataset
\newcommand{\crdataset}[1]{{\dataset{#1}^{CR}}}
\newcommand{\srdataset}[1]{{\dataset{#1}^{SR}}}

%% Detecting events
\newcommand{\detector}{\delta}            % Detector function

%% Hypothesis testing
\newcommand{\hypothesis}[1]{{\mathcal{H}_{#1}}}
\newcommand{\nullhypothesis}{{\hypothesis{0}}}
%%% Parameterizations and pdfs
\newcommand{\pdf}{{\lambda}}
\newcommand{\param}{{\mu}}
\newcommand{\nparam}{{\theta}}
\newcommand{\parametrized}[1]{{#1}_{\nparam}^{\param}}
\newcommand{\parampdf}{{\parametrized{\pdf}}}
\newcommand{\paramf}{{f^{\param}}}

%% Hypotheses
%%% Distributions
\newcommand{\dist}[1]{{\pdf_{#1}(\anyevent)}}   % Any distribution of x
\newcommand{\srdist}[1]{{\pdf_{#1}^{SR}(\anyevent)}}   % Any distribution of x
\newcommand{\crdist}[1]{{\pdf_{#1}^{CR}(\anyevent)}}   % Any distribution of x
\newcommand{\distevent}[1]{\pdf_{#1}(\event)}
\newcommand{\srdistevent}[1]{\pdf_{#1}^{SR}(\event)}
\newcommand{\crdistevent}[1]{\pdf_{#1}^{CR}(\event)}
\newcommand{\refdist}[1][\anyevent]{n_{\mathcal{R}}({#1})}  % Reference distribution
\newcommand{\crrefdist}[1][\anyevent]{n_{\mathcal{R}}^{CR}({#1})}
\newcommand{\srrefdist}[1][\anyevent]{n_{\mathcal{R}}^{SR}({#1})}

%% Likelihood
\newcommand{\likelihood}{\mathcal{L}}           % Likelihood function
\newcommand{\llikelihood}{\log\mathcal{L}}      % Log-likelihood function
\newcommand{\plratio}{\Lambda(\param)}
\newcommand{\invllratio}{{q_{0}}}

% =============================================================================
% Theorems
% =============================================================================

\newtheorem{theorem}{Theorem}
\newtheorem{definition}[theorem]{Definition}
\newtheorem{conjecture}[theorem]{Conjecture}
\newtheorem{lemma}[theorem]{Lemma}
\newtheorem{notation}[theorem]{Notation}
\newtheorem{observation}[theorem]{Observation}
\newtheorem{proposition}[theorem]{Proposition}
\newtheorem{property}[theorem]{Property}
\newtheorem{example}[theorem]{Example}
\newtheorem{remark}[theorem]{Remark}
\newtheorem{corollary}[theorem]{Corollary}

\newtheorem*{theorem*}{Theorem}

% =============================================================================
% Display Macros
% =============================================================================

% Size relative to the enclosing column, minipage, or slide column.
\newcommand{\addsizedgraph}[2]{%
        \centering
        \includegraphics[
            width=#1\linewidth,
            keepaspectratio,
            trim=0 24bp 0 0,
            clip
        ]{#2}
}
\newcommand{\addgraph}[1]{\addsizedgraph{1}{#1}}

\newcommand{\addsizedpic}[2]{%
        \centering
        \includegraphics[
            width=#1\linewidth,
            keepaspectratio
        ]{#2}
}
\newcommand{\addpic}[1]{\addsizedpic{1}{#1}}

% Visible draft replacement for figures whose assets are not available yet.
% Unlike \includegraphics{}, this keeps incomplete manuscripts compilable.
\newcommand{\placeholdergraphic}[1][Figure pending]{%
    \centering
    \fbox{\parbox[c][0.18\textheight][c]{0.85\linewidth}{%
        \centering\textit{#1}%
    }}%
}

% Macro to insert a full-page PDF without altering document orientation
\newcommand{\includefullpage}[2][1]{%
  \clearpage
  \thispagestyle{empty}% Suppress header/footer for the inserted page
  \noindent
  \begin{picture}(0,0)
    \put(-72,-750){% Adjust offset if needed for your margins
      \includegraphics[page=#1,width=\paperwidth,height=\paperheight,keepaspectratio=false]{#2}%
    }
  \end{picture}%
  \clearpage
}

% =============================================================================
% DRAFT MARKERS (TODO/FIXME/TBD)
% =============================================================================
% Usage:
%   \TODO{finish this paragraph}
%   \FIXME{replace placeholder with measured value}
%   \TBD{insert citation}
%
% Release behavior:
%   - Default is draft mode (markers visible in PDF).
%   - Set \draftfalse (after loading definitions) to make unresolved markers
%     fail the build.
\makeatletter
\@ifundefined{ifdraft}{\newif\ifdraft\drafttrue}{}
\makeatother

\ifdraft
  \newcommand{\draftmarker}[2]{%
    \textcolor{#1}{\textbf{[#2]}}%
  }
  \newcommand{\TODO}[1]{\draftmarker{red}{TODO: #1}}
  \newcommand{\FIXME}[1]{\draftmarker{orange}{FIXME: #1}}
  \newcommand{\TBD}[1]{\draftmarker{blue}{TBD: #1}}
\else
  \newcommand{\TODO}[1]{\errmessage{Unresolved TODO: #1}}
  \newcommand{\FIXME}[1]{\errmessage{Unresolved FIXME: #1}}
  \newcommand{\TBD}[1]{\errmessage{Unresolved TBD: #1}}
\fi

% =============================================================================
% END OF DEFINITIONS
% =============================================================================

% Shared abbreviations for the thesis and paper.
% Use \gls{key} in the text; the first use prints the full term and acronym,
% and later uses print the acronym only.

\newacronym{bsm}{BSM}{Beyond the Standard Model}
\newacronym{cr}{CR}{control region}
\newacronym{ddp}{DDP}{Data Directed Paradigm}
\newacronym{dof}{DoF}{degrees of freedom}
\newacronym{lfv}{LFV}{lepton flavor violation}
\newacronym{ndf}{NDF}{number density function}
\newacronym{nn}{NN}{neural network}
\newacronym{pdf}{PDF}{probability density function}
\newacronym{sm}{SM}{Standard Model}
\newacronym{sr}{SR}{signal region}
\newacronym{std}{STD}{standard deviation}
\newacronym{fcc-ee}{FCC-ee}{Future Circular Collider electron--positron collider}
\newacronym{pmssm}{pMSSM}{phenomenological minimal supersymmetric Standard Model}

\makenoidxglossaries

\newtoggle{ispaper}
\newtoggle{isthesis}

\makeatletter  % set paper size if not included in presentation
\@ifclassloaded{beamer}{
    \usepackage[
        backend=biber,
        style=numeric,
        sorting=none
    ]{biblatex}
    \addbibresource{common/bibliography/bibliography.bib}
}{
    \ifdefined\paper
        % Apply this only when compiling the main paper
        \toggletrue{ispaper}
        \togglefalse{isthesis}
    \fi

    \ifdefined\thesis
        % Apply this only when compiling the thesis
        \geometry{a4paper}
        \togglefalse{ispaper}
        \toggletrue{isthesis}
    \fi

}
\makeatother

% Package additions from REVTeX templates
\usepackage{dcolumn}% Align table columns on decimal point
% REVTeX uses a legacy label implementation. On modern kernels restore the
% hook-aware kernel version before hyperref/nameref initializes references.
\ifdefined\paperKernelLabel
    \IfFormatAtLeastTF{2023-06-01}{%
        \let\label\paperKernelLabel
        \let\ltx@label\paperKernelLabel % amsmath caches the label command
    }{}
\fi
\usepackage{hyperref}% Load after packages that redefine cross-reference commands

\ifdefined\paper
    % lineno checks the kernel parbox wrapper; REVTeX extends that wrapper.
    % Load against the kernel wrapper, then retain REVTeX's implementation.
    \makeatletter
    \let\paperParboxRestore\@parboxrestore
    \def\@parboxrestore{\@arrayparboxrestore\let\\\@normalcr}
    \usepackage[mathlines]{lineno}
    \let\@parboxrestore\paperParboxRestore
    \makeatother
    \linenumbers\relax % Commence numbering lines in the paper
\fi
\ifdefined\thesis
    \usepackage[mathlines]{lineno}
    \linenumbers\relax % Commence numbering lines in the thesis
\fi

% Initialization of tikz lib
\usetikzlibrary{fit,calc,shadows}

% Optional geometry test block from apssamp template
%\usepackage[showframe,%Uncomment any one of the following lines to test
%%scale=0.7, marginratio={1:1, 2:3}, ignoreall,% default settings
%%text={7in,10in},centering,
%%margin=1.5in,
%%total={6.5in,8.75in}, top=1.2in, left=0.9in, includefoot,
%%height=10in,a5paper,hmargin={3cm,0.8in},
%]{geometry}

}{}
\makeatother


% Introduction
% The Introduction chapter should review and summarize the current state of knowledge in your field in terms of scientific literature, raise questions, and note mistakes or possible gaps in this knowledge.

\begin{document}

\section{Introduction}

\subsection{Beyond the Standard Model Searches}

The \gls{sm} of elementary particles~\cite{Grossman_StandardModel_2023} describes fundamental particles and their interactions, but it does not provide a complete description of nature. For example, it does not explain the nature of dark matter, the origin of neutrino masses, or the matter--antimatter asymmetry of the universe. Many \gls{bsm} theories and models have been developed to address these questions. Examples include:

\begin{itemize}
    \item Supersymmetric models, which introduce a superpartner for each \gls{sm} state and can provide dark matter candidates.
    \item Seesaw models, which introduce right-handed neutrino fields to generate neutrino masses.
    \item Leptogenesis models, which connect neutrino masses and CP violation to the matter--antimatter asymmetry.
\end{itemize}

Particle accelerators are a major tool in searches for \gls{bsm} physics because they combine high collision energies with high luminosities. Examples include the Large Hadron Collider (LHC)~\cite{CERN_Courier_LHC_Run3}, with a collision energy of $13.6\,\mathrm{TeV}$ and an expected total integrated luminosity of $550\,\mathrm{fb}^{-1}$ since its first collisions, and the proposed FCC-ee~\cite{Benedikt_2026_FCC}, with energies from $90$ to $365\,\mathrm{GeV}$ and annual integrated luminosities from $0.67$ to $68\,\mathrm{ab}^{-1}$, depending on the interaction point.

Since the discovery of the Higgs boson candidate, hundreds of searches using LHC data have found no compelling evidence for new processes beyond the \gls{sm}; see, for example, Ref.~\cite{PDG_2024} and the references therein. Nevertheless, only a small fraction of possible scenarios has been probed. For example, about $40\%$ of the pMSSM parameter points excluded by ATLAS evade the available simplified-model constraints~\cite{Ambrogi_2018_pMSSM_misses}. The large, complex datasets collected at colliders therefore remain a valuable resource for further searches.

A conventional search for a specific \gls{bsm} model~\cite{ATLAS_Higgs_2012} involves several manual steps. These include deriving production mechanisms and decay modes of new particles, defining event selections to enhance the signal relative to the \gls{sm} background, quantifying systematic uncertainties, and constructing a statistical model. These steps require substantial time and effort while probing only a limited region of the possible \gls{bsm} scenarios.

In previous work~\cite{Volkovich_2022_DDP, Birman_2022_Symmetries, LastDDPPaper}, we proposed our \gls{ddp} for model-agnostic searches to address these limitations. By scanning detector data for deviations from \gls{sm} predictions, this approach can identify starting points for further investigation. In particular, symmetries of the \gls{sm} impose selection rules and relations between observable processes, providing a reference against which to search for deviations.

Here, we generalize the symmetry method first proposed in Ref.~\cite{LastDDPPaper}. That work introduced the symmetrized formalism, a statistical method for comparing pairs of data samples.

We introduce mechanisms to handle the complexities of simulated data and use them to compare sample pairs expected to be symmetric under the \gls{sm}. These extensions allow us to examine multiple observables and model nuisance parameters to account for systematic effects on the samples.

We continue our investigation of the lepton sector in search of lepton flavor violation, a well-motivated target for \gls{bsm} searches suited to this method. We use CMS Open Data sample pairs selected by lepton flavor and expected to exhibit the corresponding \gls{sm} symmetry. The following sections provide the scientific background, hypothesis-testing framework, and prior work underlying this investigation.

\subsection{Scientific Background}

The \gls{sm}~\cite{Grossman_StandardModel_2023} accounts for the strong, weak, electromagnetic and Yukawa interaction in nature. It postualtes the fields and the local symmetries imposed on the Lagrangian, and obtains the interactions, accidental symmetries and parameters that stem from its postulates.

We first formulate the total symmetry group under which the \gls{sm} is invariant:

\begin{equation}
    \label{eq:SM_symmetry_group}
    SU(3)_C \times SU(2)_L \times U(1)_Y
\end{equation}

\noindent along with tree fermion generations, each consisting of five representations:

\begin{equation}
    Q_{L_i}(3, 2)_{+1/6}, U_{R_i}(3, 1)_{+2/3}, D_{R_i}(3, 1)_{-1/3}, L_{L_i}(1, 2)_{-1/2}, E_{R_i}(1, 1)_{-1}
\end{equation}

\noindent for $i=1,2,3$. To this a scalar multiplet is added:

\begin{equation}
    \phi(1, 2)_{+1/2}
\end{equation}

\noindent while the notation used is interpreted such that both scalars in the parentheses followed by the one subscript to them are the field's representation under the according symmetry groups in Eq.~\ref{eq:SM_symmetry_group}.

This leads to spontaneous symmetry breaking. It is defined that the \gls{sm} exhibits the breaking pattern:

\begin{equation}
    \label{eq:sm_symmtery_breaking_pattern}
    SU(3)_C \times SU(2)_L \times U(1)_Y \to SU(3)_C \times U(1)_{EM}
\end{equation}

\noindent the new group being the electromagnetic group. The fermions that transform as singlets under the color group are called leptons.

Leptons appear in 3 generations, and have their masses arise from the Yukawa part of nature's Lagrangian:

\begin{equation}
    \label{eq:yukawa_lagrangian}
    -\mathcal{L}_{\text{Yuk}}
    = Y^{u}_{ij}\overline{Q_{L_i}}U_{R_j}\tilde{\phi}
    + Y^{d}_{ij}\overline{Q_{L_i}}D_{R_j}\phi
    + Y^{e}_{ij}\overline{L_{L_i}}R_{R_j}\phi
    + \text{h.c.}
\end{equation}

\noindent where $i, j = 1, 2, 3$ are the flavor indices, the $Y$ matrices are the complex $3 \times 3$ Yukawa couplings for the up, down and lepton sectors of the model, accordingly. Without loss of generality, we can choose a basis in which the $Y^{e}$ is diagonal, for which the eigenstates are the \glspl{sm} leptons. The PMNS matrix, the leptonic mixing matrix, stems from $Y^e$.

The lepton left handed fields are denoted by:

\begin{equation}
    \label{eq:lepton_lh_fields}
    L_{L_e} = \begin{pmatrix}
        \nu_{eL} \\
        e_{eL}
    \end{pmatrix}, \quad
    L_{L_\mu} = \begin{pmatrix}
        \nu_{\mu L} \\
        \mu_L
    \end{pmatrix}, \quad
    L_{L_\tau} = \begin{pmatrix}
        \nu_{\tau L} \\
        \tau_L
    \end{pmatrix},
\end{equation}

\noindent and the right handed ones:

\begin{equation}
    \label{eq:lepton_rh_fields}
    e_R, \quad \mu_R, \quad \tau_R
\end{equation}

Since the neutrino fields have no right handed partners, they remain massless. Eventually, the \gls{sm} predicts the existance of the particles:

\begin{itemize}
    \item electrons, $e^\pm$, and a corresponding massless electron neutrino, $\nu_e$.
    \item muons, $\mu^\pm$, and a corresponding massless muon neutrino, $\nu_\mu$.
    \item taus, $\tau^\pm$, and a corresponding massless tau neutrino, $\nu_\tau$.
\end{itemize}

\subsubsection{Neutrino Oscillations}

An accidental symmetry of the \gls{sm} is the lepton flavor symmetrty:

\begin{equation}
    \label{eq:lepton_flavor_symmetry}
    U(1)_e \times U(1)_\mu \times U(1)_\tau.
\end{equation}

\noindent it constrains that the number of particles of each flavor is conserved throughout a process. This being said, evidence from the Kamiokande experiment~\cite{SuperKamiokande_1998_Oscillations} showed evidence that this is not the full description of nature. Running the detector experiment for 535 days, the team exhibits discrapency between the expected and the measured number $\nu_\mu$ that pass through the detector, and that this amount varies depending on the distacnce the particles travel from the sun. A possible explanation for this are "neutrino oscillations", the term coined to express that the PMNS matrix is not diagonal in the mass basis of the leptons, causing traveling leptons to change their mass state via the free space propagator.

\subsection{Searches for Deviations from the Standard Model}

In order to estimate from an experimental sample whether or not we're observing new physics, we need to formulate a hypothesis testing mechanism.

A hypothesis~\cite{Cowan_2013_Statistics_LHC} is a statement about the probability to find some sample $x$. $P(x|H)$ expresses the probability to find $x$ given that the hypothesis $H$ is true. The likelihood of the hypothesis is then, given a sample, $\likelihood(H|x) = P(x|H)$.

Considering a null hypothesis, $\nullhypothesis$, which we want to test, and an alternative signal hypothesis, $\hypothesis{1}$. The Neymann-Pearson lemma states that the most powerful test is the (inverse) log likelihood ratio, defined as:

\begin{equation}
    \invllratio = 2 \log \dfrac{\likelihood(\hypothesis{1}|x)}{\likelihood(\nullhypothesis|x)}
\end{equation}

In our context, signal hypotheses are hypotheses spaces - parametrized by parameters of interest, $\param$. In addition, both the signal and the null hypotheses are parametrized by nuisance parameters, $\nparam$'s, that are not of interest for our hypothesis and represent system uncertainties.

In this case, the log likelihood ratio would maximize over each hypothesis space to indicate which is more likely, and becomes:

\begin{equation}
    \invllratio = 2 \log \left( \dfrac{\max\limits_{\param, \nparam} \likelihood({\param},{\nparam})}{\max\limits_{\nparam} \likelihood(\param = 0, {\nparam})} \right)
\end{equation}

\subsection{Previous Works of the Symmetrized Formalism}

The symmetrized formalism, presented in \cite{Italian_no_1_Learning_New_Physics_from_a_Machine, Italian_no_2_Learning_New_Physics_from_an_Imperfect_Machine}, suggests a method to perform an agnostic new physics search based on known measures of likelihood ratio. They do this by comparing symmetric samples rather than modeling the physics of the background and signal channels. They did so by applying neural networks to solve complex minimization expressions while maintaining a great expressiveness of the signal and null hypotheses spaces.

the goal of the method is to solve for maximization of the test statistic:

\begin{equation}
\label{eq:italian_t_definition}
    t(\mathcal{D}, \mathcal{A})
    =
    2 \log \left(
    \frac{
        \displaystyle \max_{\mathbf{w},\nu}
        \left[
            \mathcal{L}
            \left(
                \mathrm{H}_{\mathbf{w},\nu}
                \mid
                \mathcal{D},\mathcal{A}
            \right)
        \right]
    }{
        \displaystyle \max_{\nu}
        \left[
            \mathcal{L}
            \left(
                \mathrm{R}_{\nu}
                \mid
                \mathcal{D},\mathcal{A}
            \right)
        \right]
    } \right) .
\end{equation}

with $\mathbf{w}$ and $\nu$ the parameters of interest and the nuisance parameters accordingly, $\likelihood$ is the likelihood of a parametrized signal (null) hypothesis, $\mathrm{H}$ ($\mathrm{R}$), given the data sample, $\mathcal{D}$, and the auxiliary sample ,$\mathcal{A}$.

Hypotheses are formulated as:

\begin{equation}
    n(x|\mathrm{H}_{\mathbf{w}, \nu}) = e^{f(x;\mathbf{w}}n(x|\mathrm{R}_\nu)
\end{equation}

\noindent use the output of a neural network for $f$, allowing to later fit it to minimize a loss, extracted from the expressions in Equation~\ref{eq:italian_t_definition}. The exponantiated $f$ shifts a common "null" distribution, $n(x|\mathrm{R}_\nu)$, to allow for easy comparison between the two.

Our previous work \cite{LastDDPPaper} applies the symmetric formalism for samples of arbitrary sizes, removing the requirement for one of them to be significantly larger than the other, while solving several computational difficulties arose there.

with a similar test statistic objective, we redefined the hypothesis spaces as can be seen in Equation~\ref{eq:last_ddp_paper_hypotheses}. Maintaining the loss fitting mechanism for $f$ and $g$, we now have a two-fold hypothesis space, accounting for either sample's generating distribution. Both exponents modify a common reference distribution, $\refdist$, for straight-forward comparison.

\end{document}